\section{RC High Pass Filter}

\paragraph{Overview:}
A simple two component circuit that passes high frequency
signals and attenuates low  frequency signals. It is a
single pole filter but can be used in multiple stages to
create higher order filters.

\paragraph{Schematic:}
\begin{center}
  \begin{circuitikz}
    % Paths, nodes and wires:
  	\draw (9, 7.75) to[american voltage source, l={$V_{cc}$}] (9, 6.5);
  	\node[ground] at (9, 6.5){};
  	\draw (10.75, 5.5) to[european resistor, l={$G_1$}] (10.75, 4);
  	\draw (10.75, 7.5) to[capacitor, l={$C_1$}] (10.75, 6.25);
  	\draw (9, 7.75) -| (9, 8) -- (10.25, 8);
  	\draw (10.75, 5.75) -- (11.75, 5.75);
  	\node[shape=rectangle, minimum width=0.965cm, minimum height=0.465cm](N1) at (12.25, 5.75){} node[anchor=center] at (N1.text){$V_{out}$};
  	\draw (9, 3.75) to[american voltage source, l_={$V_{ee}$}] (9, 5);
  	\node[ground, xscale=-1, yscale=-1] at (9, 5){};
  	\draw (10.25, 8) -| (10.75, 7.5);
  	\draw (10.75, 6.25) |- (10.75, 5.5);
  	\draw (10.75, 4.25) -| (10.75, 3.5) -- (9, 3.5) -| (9, 3.75);
  \end{circuitikz}
\end{center}

\subsection{Transfer Function}
\raggedright
Applying KCL at $V_{out}$, we get:
\begin{equation}
  \begin{aligned}
    (V_{out} - V_{cc})sC_1 + (V_{out} - V_{ee}){G_1} &= 0 \\
    V_{out}sC_1 - V_{cc}sC_1 + V_{out}G_1 - V_{ee}G_1 &= 0
  \end{aligned}
\end{equation}

Now! A few important equations can be derived:
\begin{enumerate}
  \item Transfer Function: We rearrange KCL for Transfer Function
    in terms of \(\frac{V_{out}}{V_{cc}}\):
    \begin{align*}
      V_{out}(sC_1 + G_1) - V_{cc}sC_1 - V_{ee}G_1 &= 0 \\
      V_{out}(sC_1 + G_1) &= V_{cc}sC_1 + V_{ee}G_1 \\
      V_{out} &= \frac{V_{cc}sC_1 + V_{ee}G_1}{sC_1 + G_1} \\
      \boxed{\frac{V_{out}}{V_{cc}} = \frac{V_{cc}sC_1 + V_{ee}G_1}{V_{cc}(sC_1 + G_1)}}
    \end{align*}
  \item Poles: We can find the poles by setting the denominator
    of the transfer function to zero:
    \begin{align*}
      V_{cc}(sC_1 + G_1) &= 0 \\
      sC_1 + G_1 &= 0 \\
      s &= -\frac{G_1}{C_1} \\
      \text{Complex Pole at } &s = -\frac{1}{R_1 C_1} \\
    \end{align*}
    We can calculate the frequency corresponding to this pole:
    \begin{align*}
      s &= -2 \pi f_p \\
      -2 \pi f_p &= -\frac{1}{R_1 C_1} \\
      \boxed{f_p = \frac{1}{2 \pi R_1 C_1}}
    \end{align*}
  \item Input Impedance \(Z_{in}\):
    \begin{align*}
      Z_{in} &= \frac{1}{sC_1} + \frac{1}{G_1} \\
      \boxed{Z_{in} = \frac{1}{sC_1} + R_1}
    \end{align*}
  \item Output Impedance \(Z_{out}\):
    \begin{align*}
      Z_{out} &= \frac{1}{sC_1 + G_1} \\
      Z_{out} &= \frac{1}{sC_1 + \frac{1}{R_1}} \\
      Z_{out} &= \frac{1}{\frac{sC_1 R_1 + 1}{R_1}} \\
      Z_{out} &= \frac{R_1}{1 + sC_1 R_1}\frac{1 - sC_1 R_1}{1 - sC_1 R_1} \\
      Z_{out} &= \frac{R_1(1 - sC_1 R_1)}{1^2 - (sC_1 R_1)^2} \\
      \boxed{Z_{out} = \frac{R_1(1 - sC_1 R_1)}{1 - (sC_1 R_1)^2}}
    \end{align*}
\end{enumerate}

\subsection{Example Design}
\paragraph{Question:}
Design a high pass filter with a cutoff frequency of
\(f_p = 1kHz\) using a capacitor value of \(C_1 = 1\mu F\).
Also compute the input and output impedances at
\(f = 75Hz\).

\paragraph{Solution:}
\begin{enumerate}
  \item Calculating \(R_1\):
  \begin{align*}
    f_p &= \frac{1}{2 \pi R_1 C_1} \\
    R_1 &= \frac{1}{2 \pi f_p C_1} \\
    &= \frac{1}{2 \pi (1kHz)(1\mu F)} \\
    &= 159.15 \Omega \\
    \boxed{R_1 \approx 159 \Omega}
  \end{align*}
  \item Calculating \(Z_{in}\) at \(f = 75Hz\):
  \begin{align*}
    s &= j2 \pi f = j2 \pi (75Hz) \approx j471.24 \\
    Z_{in} &= \frac{1}{sC_1} + R_1 \\
    &= \frac{1}{j471.24 (1\mu F)} + 159.15 \Omega \\
    &= -j2122.06 + 159.15 \\
    \boxed{Z_{in} \approx 159.15 - j2122.06 \Omega}
  \end{align*}
  \item Calculating \(Z_{out}\) at \(f = 75Hz\):
  \begin{align*}
    Z_{out} &= \frac{R_1(1 - sC_1 R_1)}{1 - (sC_1 R_1)^2} \\
            &= \frac{159.15(1 - j471.24 (1\mu F) 159.15)}{1 - (j471.24 (1\mu F) 159.15)^2} \\
            &= \frac{159.15(1 - j0.075)}{1 - (-0.075^2)} \\
            &= \frac{159.15 - j11.94}{1 + 0.005625} \\
            &= \frac{159.15 - j11.94}{1.005624} \\
            \boxed{Z_{out} \approx 158.26 - j11.86 \Omega}
  \end{align*}
\end{enumerate}

\subsection{Non-Idealities}
\begin{itemize}
  \item \textbf{Component Tolerances:} Refer to \ref{sec:theory_general_tolerances}.
  \item \textbf{Parasitics:} Refer to \ref{sec:theory_general_parasitics}.
  \item \textbf{Loading Effects:} Refer to \ref{sec:theory_general_loading}.
\end{itemize}

\subsection{Applications}
\begin{itemize}
  \item \textbf{AC Coupling (DC Blocking):}
    Removes DC offsets while passing the AC signal.
    \textit{Used between amplifier stages to prevent bias shift.}

  \item \textbf{Audio Pop and Drift Removal:}
    Eliminates slow voltage drift and power-on transients that
    cause clicks or speaker cone offset.

  \item \textbf{Sensor Offset Removal:}
    Blocks slowly varying offsets in sensors while preserving
    fast changes.
    \textit{Example: vibration or motion detection.}

  \item \textbf{Edge and Transient Detection:}
    Converts sudden changes in voltage into sharp pulses.
    \textit{Used in timing, triggering, and reset circuits.}

  \item \textbf{Baseline Restoration:}
    Keeps signal centered around zero by removing accumulated
    low-frequency components.
\end{itemize}

